\documentclass[11pt]{article}
\usepackage[margin=0.8in]{geometry}
\usepackage{amsmath,amssymb}
\usepackage[hidelinks]{hyperref}
\setlength{\parindent}{0pt}
\setlength{\parskip}{5pt}
\newcommand{\relu}{\operatorname{ReLU}}
\begin{document}
\begin{center}
{\Large\bfseries Question 2: Manual Backpropagation}\par
{\large One iteration and verification with PyTorch}
\end{center}

\textbf{Scope.} Calculate the gradients at the notebook's initial weights,
using one input--target pair and one backward pass, before any weight update.
Vectors are row vectors; $(W_\ell)_{ij}$ connects input $i$ to output $j$.
Both layers have no bias.

\section*{1. Initial values and forward pass}
The notebook specifies
\[
X=\begin{bmatrix}1&2\end{bmatrix},\quad
T=\begin{bmatrix}0&0\end{bmatrix},\quad
W_1=\begin{bmatrix}2&1\\0&1\end{bmatrix},\quad
W_2=\begin{bmatrix}1&2\\1&2\end{bmatrix}.
\]
The hidden-layer preactivation and activation are
\begin{align*}
S_2=XW_1
 &=\begin{bmatrix}1(2)+2(0)&1(1)+2(1)\end{bmatrix}
 =\begin{bmatrix}2&3\end{bmatrix},\\
Y_2=\relu(S_2)&=\begin{bmatrix}2&3\end{bmatrix}.
\end{align*}
The output-layer preactivation and activation are
\begin{align*}
S_3=Y_2W_2
 &=\begin{bmatrix}2(1)+3(1)&2(2)+3(2)\end{bmatrix}
 =\begin{bmatrix}5&10\end{bmatrix},\\
Y_3=\relu(S_3)&=\begin{bmatrix}5&10\end{bmatrix}.
\end{align*}
All four preactivations are positive. Hence every ReLU derivative used in
this iteration is $1$.

\section*{2. Mean squared error and output error signal}
The default \texttt{nn.MSELoss()} averages over the two output elements
\cite{mse}. Therefore
\[
L=\frac{(y_{3,1}-t_1)^2+(y_{3,2}-t_2)^2}{2}
 =\frac{5^2+10^2}{2}=62.5.
\]
For each output $j$,
\[
\frac{\partial L}{\partial y_{3,j}}
 =\frac{2}{2}(y_{3,j}-t_j),\qquad
\delta_3=\frac{\partial L}{\partial S_3}
 =\begin{bmatrix}5&10\end{bmatrix}\odot\begin{bmatrix}1&1\end{bmatrix}
 =\begin{bmatrix}5&10\end{bmatrix}.
\]
Here $\odot$ denotes elementwise multiplication.

\section*{3. Gradient of the second-layer weights}
Since $s_{3,j}=\sum_i y_{2,i}(W_2)_{ij}$, the chain rule gives
\[
\frac{\partial L}{\partial(W_2)_{ij}}=y_{2,i}\delta_{3,j}.
\]
Thus all four derivatives are the entries of the outer product
\[
\boxed{\nabla_{W_2}L=Y_2^{\mathsf T}\delta_3
 =\begin{bmatrix}2\cdot5&2\cdot10\\3\cdot5&3\cdot10\end{bmatrix}
 =\begin{bmatrix}10&20\\15&30\end{bmatrix}.}
\]

\newpage
\section*{4. Hidden error signal and gradient of the first-layer weights}
Each hidden activation affects both outputs. Applying the chain rule,
\[
\frac{\partial L}{\partial y_{2,i}}
 =\sum_j\delta_{3,j}(W_2)_{ij},\qquad
\frac{\partial L}{\partial Y_2}=\delta_3 W_2^{\mathsf T}.
\]
In numbers,
\[
\delta_3 W_2^{\mathsf T}
 =\begin{bmatrix}5&10\end{bmatrix}
  \begin{bmatrix}1&1\\2&2\end{bmatrix}
 =\begin{bmatrix}5(1)+10(2)&5(1)+10(2)\end{bmatrix}
 =\begin{bmatrix}25&25\end{bmatrix}.
\]
Both hidden ReLUs have derivative $1$, so
\[
\delta_2=\frac{\partial L}{\partial S_2}
 =(\delta_3W_2^{\mathsf T})\odot\relu'(S_2)
 =\begin{bmatrix}25&25\end{bmatrix}.
\]
Since $s_{2,j}=\sum_i x_i(W_1)_{ij}$,
\[
\frac{\partial L}{\partial(W_1)_{ij}}=x_i\delta_{2,j},\qquad
\boxed{\nabla_{W_1}L=X^{\mathsf T}\delta_2
 =\begin{bmatrix}1\cdot25&1\cdot25\\2\cdot25&2\cdot25\end{bmatrix}
 =\begin{bmatrix}25&25\\50&50\end{bmatrix}.}
\]

\section*{5. Verification and the \texttt{nn.Linear} transpose}
The notebook's direct \texttt{torch.matmul} implementation prints
\[
\texttt{W1.grad}=\begin{bmatrix}25&25\\50&50\end{bmatrix},\qquad
\texttt{W2.grad}=\begin{bmatrix}10&20\\15&30\end{bmatrix}.
\]
These exactly match the manual derivatives above.

PyTorch's \texttt{nn.Linear} computes $xA^{\mathsf T}+b$ and stores $A$ with
shape (output features, input features) \cite{linear}. The notebook sets
\[
A_1=W_1^{\mathsf T}=\begin{bmatrix}2&0\\1&1\end{bmatrix},\qquad
A_2=W_2^{\mathsf T}=\begin{bmatrix}1&1\\2&2\end{bmatrix}.
\]
Thus $XA_1^{\mathsf T}=XW_1$ and $Y_2A_2^{\mathsf T}=Y_2W_2$.
Since $(A_\ell)_{ji}=(W_\ell)_{ij}$, their parameter gradients obey
\[
\nabla_{A_\ell}L=(\nabla_{W_\ell}L)^{\mathsf T}.
\]
The executed notebook confirms
\[
\texttt{fc1.weight.grad}=\begin{bmatrix}25&50\\25&50\end{bmatrix},\qquad
\texttt{fc2.weight.grad}=\begin{bmatrix}10&15\\20&30\end{bmatrix}.
\]
Comparing the notebook's first backward-pass outputs with the manual matrices
gives maximum absolute error $0$ for both direct gradients and both transposed
\texttt{nn.Linear} gradients. Both implementations give
$Y_3=[5,10]$ and $L=62.5$ before the first update.

\begingroup\small
\begin{thebibliography}{2}
\bibitem{linear} PyTorch documentation, \texttt{torch.nn.Linear}.
\url{https://docs.pytorch.org/docs/2.14/generated/torch.nn.Linear.html}
\bibitem{mse} PyTorch documentation, \texttt{torch.nn.MSELoss}.
\url{https://docs.pytorch.org/docs/2.14/generated/torch.nn.MSELoss.html}
\end{thebibliography}
\endgroup
\end{document}
